\documentclass[../main.tex]{subfiles}
\begin{document}

\rule{0pt}{\textheight}
\newpage
\rule{0pt}{\textheight}
\parbox[b]{0.85\textwidth}{
    \centering
    \emph{El formato de este libro está basado en The Legrand Orange Book, por Mathias Legrand, y todo el contenido está bajo licencia Creative Commons Atribución - No comercial - Compartir Igual - No portada 3.0(CC BY-NC-SA 3.0), disponible en \emph{\url{https://creativecommons.org/licenses/by-nc-sa/3.0/deed.es}}}
       \includegraphics[scale=1]{cc/cc.pdf} 
       \includegraphics[scale=1]{cc/by.pdf} 
       \includegraphics[scale=1]{cc/nc.pdf} 
       \includegraphics[scale=1]{cc/sa.pdf} 
    \vspace{0.5\textheight}
    \vfill 

    \itemizePlainStyle
    \itshape
    \raggedright
    Autores:
    \begin{itemize}
        \item Mariana Valoy
        \item Ana Merino
        \item Laura González 
        \item Ágata Ordano
        \item Cecilia Ordano
    \end{itemize}\hfill\\[3cm]
    \normalfont
    \restoreItemizeStyle
    \centering

    \vfill 
    \emph{Este manual fue escrito en diciembre de 2020 para la Fundación Nodo Ambiental -- \emph{\url{https://nodoambiental.org}}, y puede no ser representativo de cambios ocurridos después de la fecha.}
}
\rule{0pt}{\textheight}

\thispagestyle{empty}
\pagebreak

\tableofcontents

\chapter*{Cómo usar este manual}

Este manual está dividido en tres partes:

\begin{itemize}
    \item Capítulo 1 -- La huerta agroecológica y sus fundamentos.
    \item Capitulo 2 -- Manejo de fauna y biodiversidad asociadas a la huerta.
    \item \color{black!25}{Capítulo 3 -- El compostaje.} \observation{Este capítulo no existe}
\end{itemize}

\hfill\\

El contenido de este libro es asistido por marcadores de colores especiales que representan conceptos determinados:

\begin{itemize}
    \item \tcbox{Elemento clave}: Este marcador representa un \emph{elemento clave}, y lo que contenga es un concepto importante.
    \item recuadro lima:\begin{recuadroV}
            Representa un concepto o grupos de conceptos que son significativos.
    \end{recuadroV}
    \item recuadro salmón:\begin{recuadroR}
            Representa un concepto o grupos de conceptos que requieren atención especial de alguna manera, siendo posibles fuentes de problemas.
    \end{recuadroR}
\end{itemize}

\vfill

Además, este documento está escrito como una guía, de manera tal que se puede consultar cualquier parte del mismo como referencia; no es necesario leer esto en orden cronológico, aunque se recomienda hacerlo.




\end{document}